\section{Continuous Monitoring Pipeline}

\label{sec:monitoring}
% ══════════════════════════════════════════════════════════════════════════════

\subsection{Data Ingestion}

ZooMRV ingests data from multiple satellite constellations:

\begin{table}[H]
\centering
\caption{Satellite data sources and processing characteristics.}
\label{tab:satellites}
\begin{tabular}{p{3cm}p{2cm}p{2cm}p{2cm}p{2.5cm}}
\toprule
\textbf{Source} & \textbf{Resolution} & \textbf{Revisit} &
  \textbf{Coverage} & \textbf{Primary use} \\
\midrule
Sentinel-2 & 10\,m & 5 days & Global & AGB, change detection \\
Sentinel-1 SAR & 10\,m & 6 days & Global & Cloud-free monitoring \\
Landsat-9 & 30\,m & 16 days & Global & Long-term trends \\
GEDI LiDAR & 25\,m & --- & Tropical & Canopy height, AGB cal. \\
Planet NICFI & 4.7\,m & Monthly & Tropical & Fine-scale change \\
\bottomrule
\end{tabular}
\end{table}

\subsection{Change Detection}

We implement a multi-scale change detection algorithm:

\paragraph{Pixel-level.} A temporal convolutional network (TCN) processes
per-pixel time series of NDVI, NBR, and SWIR2 bands to detect abrupt
changes (logging, fire) and gradual degradation:
\begin{equation}
  \hat{c}_i = \text{TCN}(\mathbf{x}_{i,t-T:t}) \in [0,1]
\end{equation}
where $\hat{c}_i$ is the change probability for pixel $i$.

\paragraph{Object-level.} A U-Net semantic segmentation model
identifies changed areas at the patch level, providing spatial context
that reduces false positives from atmospheric and phenological variation.

\paragraph{Alert aggregation.} Pixel and object-level detections are
aggregated into alert polygons with associated confidence scores,
change type classification (deforestation, degradation, fire, flood),
and estimated biomass loss.

\subsection{Reversal Detection Performance}

\begin{table}[H]
\centering
\caption{Reversal detection performance across ecosystem types.}
\label{tab:reversals}
\begin{tabular}{lccc}
\toprule
\textbf{Ecosystem} & \textbf{Detection rate} &
  \textbf{False positive rate} & \textbf{Median latency} \\
\midrule
Tropical forest & 94.2\% & 3.1\% & 8 days \\
Mangrove & 91.7\% & 4.8\% & 11 days \\
Peatland & 87.3\% & 6.2\% & 18 days \\
Grassland/savanna & 83.6\% & 8.1\% & 15 days \\
\midrule
\textbf{Overall} & \textbf{91.4\%} & \textbf{4.7\%} &
  \textbf{12 days} \\
\bottomrule
\end{tabular}
\end{table}

ZooMRV detects 91.4\% of reversal events with a median latency of 12
days, compared to the 14-month average for traditional audit cycles.
The system detects 23\% more reversals than periodic audits, primarily
small-scale degradation events that fall below audit detection
thresholds.


% ══════════════════════════════════════════════════════════════════════════════
